\section{Diskussion}
\label{sec:Diskussion}
Zur Berechnung der relativen Abweichung von experimentell bestimmten Werten zu theoretischen Werten wird 
\begin{equation*}
    \increment x_{\mathrm{rel}} = \left| \frac{x_{\mathrm{exp}}-x_{\mathrm{theo}}}{x_{\mathrm{theo}}}\right|
\end{equation*}
verwendet. 
\\
Die Ergebnisse der Untersuchung der transversalen Moden in \autoref{subsec:tem} sind zufriedenstellend, da jeweils der erwarteten Intensitätsverteilungen deutlich zu erkennen sind.
Das Abweichen einiger Werte lässt sich durch verschiedene Faktoren erklären. 
Da die Messung bei angeschaltetem Raumlicht durchgeführt wird, könnten sich Fehler durch Schwankungen in dessen Intensität ergeben. 
Außerdem kann die verwendete Photodiode für Schwankunen gesorgt haben.
So fiel bei der Durchführung bereits auf, dass hintere Nachkommastellen stark geschwankt haben.
Zwar wurden diese abgeschnitten und sind in Messunsicherheiten eingeflosse, aber es ist trotzdem denkbar, dass auch einige der Werte durch solche Schwankunden verfälscht wurden. 
Zudem können außere elektromagnetisch Felder Einfluss auf das Messgerät genommen und so die Werte verfälscht haben. 
Darüber hinaus ist ein wichtiger Faktor, dass bei der händischen Justierung einige Elemente des Aufbaus, nicht genau gerade in den Strahlengang eingebracht wurden oder die tatsächliche optische Achse nicht genau grade durch den Aufbau verlaufen ist
Dadurch können sich durch zusätzliche Reflexionen Fehler ergeben.
Besonder bei der TEM01-Mode fällt auf, dass einer der Peaks deutlich größer ist, als der andere. 
Dies lässt sich insbesondere durch die zuletzt genannten Punkte erklären. 
Hier kann aber auch vor allem eine Schiefstellung des zur Unterdrückung der Grundmode eingebrachten Drahtes eine Rolle gespielt haben. 
So war dessen Einbringen in den Resonator schweirig, da sich der Draht leicht schief in seiner Halterung befand.
So kann sich ungleiche Höher der Peaks des gemessenen Intensitätsprofils auch durch Schattenwurf des Drahtes ergeben haben.
\\
Auch bei der Messung der Stabilitätsbedinung gibt es leichte Abweichungen in den Messwerten, obwohl der theoretische Verlauf deutlich zu erkennen ist.  
Für die maximale Resonatorlänge bei Anordnung mit Planspiegel und konkavem Spiegel ergibt sich eine relative Abweichung von $\increment L_{\mathrm{max}}=25,07 \%$ zum theoretischen Wert.
Diese Abweichungen lassen sich vor allem durch menschliche Fehler beim Justieren erklären, da hier für jeden Messung der Laser neu einjustiert werden muss. 
Dabei können sich, wie beim der Untersuchung der transversalen Moden, jedes mal neue Schiefstellungen der optischen Achse und der Bauelemente ergeben haben, die die Ergebnisse verfälschen. 
So ist auch zu erklären, das für eine größere Resonatorlänge, keine Laseraktivitär mehr erzeugt wurde. 
Zudem können auch hier Fehler in dem Messinstrument für die Intensität vorliegen. 
So lässt sich auch der negative Dunkelstrom, durch eine fehlerafte Einstellung des Messinstrumentes erklären. 
Auch hier kann das Raumlicht eine Rolle bei den Fehlern gespielt haben.
\\
Die Messwerte der Polratisation passen sehr gut zu dem erwarteten theoretischen Verlauf. 
Leichte Abweichungen einzelner Messwerte lassen sich durch die bereits zu den anderen beiden Messungen genannten Punkte erklären. 

